\begin{textsample}{POS Dim 5 – global\_north\_2025\_04 – Score 20.00 – global\_north\_2025\_04/124128095.txt}  \label{ex:f5_pos_091}
The only certainty that the Palestinian people have is that tomorrow is going to be worse said the South \textbf{African} representative at the ICJ hearing.
A Palestinian demonstrator holds a sign thanking South Africa for its support during a protest in Amman, Jordan. ( photo credit : JEHAD SHELBAK/REUTERS ) South Africa called for Israel to follow the \textbf{provisional} \textbf{measures} issued to it by the \textbf{International} \textbf{Court} of Justice and to reverse its decision to ban the United Nations Relief and Works Agency."The world has failed the Palestinian people.
The only certainty they have is that tomorrow is going to be worse,"said the South \textbf{African} representative.
South Africa is one of 40 countries delivering remarks before the ICJ this week after the \textbf{court} was tasked in December with forming an \textbf{advisory} opinion on Israel ' s \textbf{obligations} to facilitate aid to Palestinians, delivered by states and \textbf{international} groups, including the United Nations. \textbf{Advisory} opinions of the ICJ carry \textbf{legal} and political weight, although they are not binding, and the \textbf{court} @ @ @ @ @ @ @ @ @ @ World \textbf{Court} will likely take several months to form its opinion.
The hearings this week will focus chiefly on the issue of starvation and aid and discuss whether Israel -- a signatory to the UN Charter -- \textbf{acted} against its commitments by banning UNRWA in November.
South Africa said this led to a lack of food and services, as did Israel ' s March ban of aid, which piled \textbf{international} pressure on the case.South \textbf{African} Ambassador to the Netherlands Vusimuzi Madonsela uses a phone at the ICJ during a \textbf{ruling} on South Africa ' s request to order a halt to Israel ' s Rafah offensive in Gaza as part of a larger \textbf{case} brought before the Hague-based \textbf{court} by South Africa accusing Israel of \textbf{genocide}, ( credit : JOHANNA GERON/REUTERS ) Israel has said it would not allow the entry of goods and supplies into Gaza until Hamas releases all 59 remaining hostages and that the food and aid that came into the enclave this year is enough to last for about another month.
Palestinian representatives testifying before the \textbf{court} on Monday emphasized @ @ @ @ @ @ @ @ @ @ a lack of food and medicine, worsened through UNRWA ' s ban and the halting of aid.
Israel has said that during the ceasefire this year between January 19 and March 18, 600 trucks of aid passed daily into the Gaza Strip, providing Gazans with enough food to last an estimated three to six months.
The United States will deliver remarks on Wednesday, along with Jordan and Iran, while Qatar, a mediator in the hostage and ceasefire negotiations alongside Egypt, will deliver remarks on Thursday.
South Africa filed the larger \textbf{genocide} \textbf{case} against Israel for alleged \textbf{violations} of the 1948 \textbf{Convention} on the \textbf{Prevention} and Punishment of the \textbf{Crime} of \textbf{Genocide} in December 2023.
Hearings took place in January 2024.
The country argued that Israel committed \textbf{genocide} by failing to stop or prosecute public \textbf{incitement} against the sentiment and through its military and aid policy in Gaza.
ICJ issued a \textbf{provisional} order After two weeks of debates, in January 2024, the ICJ issued a \textbf{provisional} order that Israel must refrain @ @ @ @ @ @ @ @ @ @ short of demanding that it stop the war.
The \textbf{court} was critical of Israel ' s actions in Gaza but made do with calling on Jerusalem to prevent any future possible \textbf{genocidal} \textbf{acts} and to ensure the immediate entry of aid.
It also called on Israel to punish those who express \textbf{incitement} against Palestinians.
This decision was approved by a large majority of the ICJ panel, 15 to two.
Israel has denied these claims, insisting that it sticks to \textbf{international} \textbf{law} and is not using starvation as a \textbf{method} of warfare to enable \textbf{genocide} against the Palestinian people but rather is fighting Hamas, a terrorist organization deeply embedded in civilian infrastructure in the Gaza Strip.

% matched lemmas: act, advisory, african, case, convention, court, crime, genocidal, genocide, incitement, international, law, legal, measure, method, obligation, prevention, provisional, ruling, violation
\end{textsample}
